\documentclass[12pt,a4paper]{article}
\usepackage[utf8]{inputenc}
\usepackage[T1]{fontenc}
\usepackage[english]{babel}
\usepackage[margin=2.5cm]{geometry}
\usepackage{url}

\setlength{\parindent}{0pt}
\setlength{\parskip}{3pt}
\linespread{1.05}

\newcommand{\frage}[1]{\par\vspace{0.7em}\noindent\textbf{#1}\par}
\newcommand{\antwort}[1]{\par\noindent\textit{Answer:} #1\par}

\begin{document}

\begin{center}
  {\LARGE\bfseries Leader's Guide with Answers}\\[4pt]
  {\large Mark 15:20--41 · Home group, 8 October 2026 · for the leader only}\\[4pt]
  {\footnotesize The numbering corresponds to the questions in the main document.}
\end{center}

\section*{Getting Started}

\subsection*{Language}

\frage{1. Which terms in the text are unclear or unfamiliar? Collect them.}
\antwort{Possible collection with brief explanations: \textbf{purple robe} -- a soldier's garment in imperial colour, used as a mocking prop (an ``emperor's robe''); \textbf{wine mixed with myrrh} -- a numbing painkiller (cf.~Prov 31:6); \textbf{Golgotha} -- Aramaic for ``skull (place)'', probably named after the shape of the hill; \textbf{third/sixth/ninth hour} -- Jewish time-keeping from 6~a.m., i.e.~9 a.m., noon, 3 p.m.; \textbf{Eloi, Eloi, lama sabachthani} -- Aramaic; Mark translates it in the text; \textbf{Elijah} -- the prophet whose return was expected at the end of days (Mal 3:23); \textbf{veil} -- the great curtain before the Most Holy Place in the temple, the symbol of the separation between God and humanity. For the discussion it matters that the group collects the word-questions first -- they resolve many problems of understanding before the actual analysis begins.}

\subsection*{Place \& Time}

\frage{2. Why is the place called ``the place of a skull'', and why was it outside the city (v.~22)?}
\antwort{\glqq Golgotha\grqq\ is Aramaic \emph{gulgalta} = skull; the translation stands in the text itself (likewise Mt 27:33). \textbf{Three explanations for the name:} (1) Most likely: the rocky hill outside Jerusalem's gates looked from a distance like a skull or a bald head -- hence the name as an ordinary place-name, not a theological construct. (2) A folkloristic explanation from later centuries: the skulls of executed people lay there -- unlikely, because Jewish purity law (contact with bones defiles) ruled out the city leading past a place of bones. (3) According to John 19:41 there was a garden with a new tomb nearby; a skull as stone or garden ornament is conceivable but remains speculation. \textbf{The textual finding matters:} Mark translates the name -- the readership hears a death-symbol at the first mention, without the text explaining any symbolism. \textbf{On the location outside the city:} executions took place before the gates -- for three converging reasons: (a) \textbf{Jewish purity law} (the decisive reason in Jerusalem): contact with a corpse defiles (Num 19:11--16), a hanged man must not remain overnight (Deut 21:22f); the Torah brought delinquents ``outside the camp'' (Num 15:35; Lev 24:14) -- Jerusalem lived this as ``outside the gate''. (b) \textbf{Roman practice}: crucifixions stood along roads and at city limits -- as deterrence (visible for days to everyone entering or leaving; Josephus reports hundreds of crucifixions daily within sight of the city in the year 70), because the punishment as staged spectacle needed an audience (so too v.~29: ``those who passed by''), and for practical reasons (cries, birds, decay did not belong in the city). (c) \textbf{Social degradation}: being led out marked social death -- the executed was removed from the community like one accursed. Heb 13:12f takes this up and reverses it: ``Jesus suffered outside the gate in order to sanctify the people -- let us then go to him outside the camp.'' For Mark 15 this means: the ``leading out'' (vv.~20.22) is part of the punishment itself -- and theologically charged: Jesus dies outside, at the edge, in the unclean. The place was also publicly accessible: passers-by can read the inscription, mock, stand and stay.}

\frage{3. Collect the time references (vv.~25.~33.~34) -- how does the timeline run? Which festival is in the background?}
\antwort{Third hour (= 9~a.m.): crucifixion (v.~25). Sixth to ninth hour (= noon to 3~p.m.): darkness (v.~33). Ninth hour (= 3~p.m.): cry of death and death (vv.~34--37). The darkness encloses exactly the midday hours. The background is the Passover festival: Jerusalem is full of pilgrims, the mood politically charged -- hence the Barabbas scene (15:6--15) shortly before.}

\subsection*{Context}

\frage{4. What happened immediately before (Mk 14:43--15:19)?}
\antwort{Arrest in Gethsemane (14:43--52; Judas betrays with a kiss, the disciples flee), hearing before the high council with Jesus' confession ``I am'' and the verdict ``guilty of death'' (14:53--65), Peter's denial (14:66--72), trial before Pilate including the Barabbas release (15:1--15), flogging and mockery by the soldiers with purple robe and crown of thorns (15:16--19). Verse 20 closes this scene: ``when they had mocked him \ldots''}

\frage{5. ``Son of God'': Mk 1:1 $\leftrightarrow$ v.~39 -- who recognises it at the end, and why is that surprising?}
\antwort{Mark 1:1 titles Jesus ``Son of God''; afterwards God confirms it (baptism 1:11; transfiguration 9:7) and demons -- but no human being understands it. Only at 15:39 does a \textbf{human being} confess this title: a Roman officer, a Gentile, himself involved in the execution -- and precisely at the sight of the dying man, not of a miracle. In Mark, sonship of God is revealed not in power but at the cross.}

\subsection*{Characters}

\frage{6. Who acts, who speaks, who is silent in vv.~20--41? (When does Jesus himself speak?)}
\antwort{Acting: the soldiers (crucify, cast lots for the garments), Simon of Cyrene (unwillingly carries the cross), those passing by, the chief priests and scribes and those crucified with him (all mock), a bystander with the vinegar sponge, the Roman centurion (confesses), the women (watch from afar). \textbf{Jesus himself speaks only once}: the psalm cry in v.~34. Throughout the whole section he is the object -- everything happens to him, he hardly acts. That is narrative intent: the Sufferer, not the Fighter, stands at the centre.}

\frage{7. Who is missing -- and why? (The disciples, Mk 14:50!)}
\antwort{The disciples. They fled at the arrest (14:50), Peter denied him (14:66--72). At the cross there are only opponents, bystanders and outsiders. The absence is part of the story: in the end it is precisely strangers (the centurion, the women) who recognise what the close followers did not.}

\section*{The Crucifixion (vv.~20--28)}

\frage{8. Simon is \textbf{compelled} to carry the cross (v.~21) -- what does that mean for Mk 8:34 (``take up his cross'')?}
\antwort{Simon carries the cross unwillingly, pressed into it -- the exact opposite of the discipleship word of Mk 8:34, where the cross is \emph{voluntarily} ``taken up''. Simon is the prototype of involuntary cross-bearing. And yet: he walks the way \emph{with} Jesus. That his son Rufus is known to the community (Rom 16:13) suggests that the compelled one became a Christian. Discipleship often begins unwillingly -- but it begins.}

\frage{9. Why does Jesus refuse the wine mixed with myrrh (v.~23)?}
\antwort{The myrrh-wine is a numbing painkiller (cf.~Prov 31:6), probably a gesture of mercy from devout Jerusalem women. Jesus refuses: he wants to endure the agony consciously -- including the cry of God-forsakenness (v.~34). Whoever is numbed cannot lament. The way of suffering is not to be anaesthetised but to be passed through.}

\frage{10. ``King of the Jews'' (v.~26) -- written as a charge: in what sense is it true?}
\antwort{The titulus named the offence; ``King of the Jews'' is meant as a political charge (insurrection against Rome). But the charge tells the truth: Jesus \emph{is} the king -- only not in the way power expects. The title occurs six times in chapter 15 (15:2.9.12.18.26.32): first from Pilate, then mockingly, finally as the inscription -- Mark turns the mockery into a Christology. The king is enthroned on the cross.}

\frage{11. \emph{Our lives}: Simon carries the cross for Jesus -- unwillingly. What does that mean for our picture of discipleship: does it begin with us -- or is it demanded of us, and then grows?}
\antwort{An important tension: Mk 8:34 demands to ``take up'' the cross -- but the text first shows someone upon whom it is laid. In Mark, discipleship always begins as a calling that comes over the disciples (``followed'', 2:14; 10:52), not as their own initiative. Peter will later himself be ``girded and led where you do not wish'' (John 21:18). Three guiding thoughts for the conversation: (a) Being carried is not the opposite of discipleship but its beginning -- Simon walks the way \emph{with} Jesus. (b) Not every suffering is automatically a ``cross'': the analogy only holds where what is endured stands in relation to Jesus. (c) The unwilling can become willing -- his son Rufus belongs to the community (Rom 16:13). Stay honest: not everything demanded later feels like a way -- but the question invites looking back, it does not force.}

\section*{Mocked on the Cross (vv.~29--32)}

\frage{12. Three groups mock -- what does each throw at him, where does something repeat itself?}
\antwort{Those passing by: temple destruction/rebuilding in three days (vv.~29f) -- an attack on his reputation as prophet. Chief priests and scribes: ``He saved others; he cannot save himself'' (v.~31) -- an attack on his saving power. Those crucified with him: ``The Christ, the King of Israel'' (v.~32) -- an attack on his messianic claim. What repeats itself: all three demand a visible sign of power (``come down'') as proof, and all three name titles of Jesus -- in the form of mockery.}

\frage{13. \textbf{The paradox of v.~31}: ``He saved others -- he cannot save himself.'' How can both be true at once? Is not-being-able-to-save-himself weakness -- or precisely the way he saves others?}
\antwort{The mockery thinks in human logic: whoever has power uses it for themselves first -- not being able to save oneself proves that the help for others was worthless. The cross reverses this logic: Jesus helps others \emph{precisely by} not saving himself -- the giving of his life is the help (Mk 10:45: ``to give his life as a ransom for many''). The paradox is not a gap in logic but the core of the message: salvation happens not \emph{despite} the cross but \emph{through} it. The chief priests involuntarily say what they do not mean: exactly in this not-saving-himself the centurion recognises the Son of God shortly afterwards (v.~39). For today: Christian help, too, often happens not from superiority but from self-giving -- and is therefore mistaken for weakness by the logic of the world.}

\frage{14. Where does the mockery become unintentionally true (vv.~29.~31.~32)?}
\antwort{In all three cases -- only differently than the mockers mean: (a) He ``destroys the temple and builds it in three days'': his death and his resurrection (so interpreted in John 2:19--21). (b) He ``saved others'': precisely by not saving himself -- the cross becomes the salvation of the many. (c) He \emph{is} the Christ and King of Israel -- as the Crucified. Mark lets the opponents proclaim the truth involuntarily: the readers recognise what the mockers do not suspect.}

\frage{15. ``\ldots\ that we may see and believe'' (v.~32) -- why is that not a genuine condition of faith?}
\antwort{Because faith that demands a visible miracle is no longer trust but compelled knowledge. The cross deliberately refuses the sign: the faith Mark seeks arises at the crucified Jesus himself -- without guarantee (cf.~John 20:29: ``Blessed are those who have not seen and yet believe''). Only after Easter does the cross become readable as a sign -- but it never becomes a proof.}

\frage{16. \emph{Our lives}: Where do we (or others) demand proof first -- and what would it mean to believe without proof?}
\antwort{Open question. Possible impulses: we live relationships without proofs -- trust grows in actions, not in demonstrations. Where do we demand signs from God (answer, turnaround) before we trust? The cross refuses the proof and gives community anyway. Stay gentle in the group: not pronouncing guilt, but speaking honestly about the need for security.}

\section*{Darkness, Cry and Confession (vv.~33--41)}

\frage{17. ``Why have you forsaken me?'' (v.~34) -- mere despair? (Psalm 22: a lament that turns into praise.)}
\antwort{Both. First, real God-forsakenness: Jesus endures the distance from the Father -- interpreted in the New Testament as bearing sin (2 Cor 5:21; Gal 3:13). Second, a word of faith: the cry quotes Ps 22:2 -- and this psalm turns into praise and proclamation among the nations (Ps 22:23--32). Whoever knows the psalm hears trust in the cry as well: Jesus addresses God despite everything as ``my God''. Lament and faith are not mutually exclusive here.}

\frage{18. Darkness and the torn veil (vv.~33.~38) -- who acts here, what is opened or judged?}
\antwort{God himself -- both signs happen without Jesus' doing. The darkness over the land (noon to 3 p.m.) is a sign of judgement and theophany (Amos 8:9; Ex 10:21f). The veil before the Most Holy Place tears ``from top to bottom'' -- from God. To be read doubly: (a) access to God is opened for all (Heb 10:19f: ``through the veil''); (b) judgement on the temple cult -- its time is over (cf.~Mk 13:2). Both together: the old separation ends, God begins somewhere else.}

\frage{19. Why does the Roman centurion of all people confess ``Son of God'' (v.~39)?}
\antwort{He is a Gentile, an officer of the occupying power, responsible for carrying out the execution -- the last person one would expect it from. His occasion is not a miracle but \emph{how} Jesus dies: ``seeing that he cried out like this and breathed his last''. The dying man himself convinces him. Thus Mk 1:1 is fulfilled: at the cross a human being recognises who Jesus is. That it is an enemy sharpens the point.}

\frage{20. Why do the women stay and see everything (vv.~40f) -- unlike the disciples?}
\antwort{The women stand ``from afar'' -- at a distance, but they \emph{stay}, while the disciples fled (14:50). They had already followed Jesus in Galilee and served him (v.~41) -- Mark names them in discipleship terminology. Their staying has a practical-theological reason: they are the only witnesses of death \emph{and} burial (v.~47) -- and therefore the credible witnesses of the empty tomb (16:1--8). The Easter message of the Gospel rests on their testimony.}

\frage{21. \emph{Our lives}: When have we experienced God's nearness or forsakenness first-hand? What helped in the darkness?}
\antwort{Open question -- here there is room for personal stories. Important: lament is biblical (Ps 22, Job, Jeremiah) and not a deficit of faith. What helped: people who ``stayed from afar''; rituals (prayer, song, communion); the knowledge of the Psalm-22 arc (lament does not end in the abyss). Whoever does not want to share their own experience may stay silent -- the question invites, it does not force.}

\frage{22. \emph{Our lives}: For whom can we today be there ``from a distance'' -- staying, watching, bearing witness?}
\antwort{Open question. Possible impulses: people in hospital, in isolation, in dying; neighbours in loneliness; also here in the home group: who is currently ``at a distance'' (hurt, moved away, closed off)? The women prove: distance is no obstacle to witnessing -- being there suffices. Transition to the prayer requests (7~p.m.).}

\section*{Cross References to Other Bible Passages}

\frage{23. Read Psalm 22 -- which verses of our text echo it?}
\antwort{Ps 22:2 $\to$ v.~34 (``My God, my God, why have you forsaken me?''); Ps 22:8 $\to$ vv.~29--32 (``He trusts in the LORD; let him deliver him'' -- the mockery of those passing by sounds like this); Ps 22:19 $\to$ v.~24 (``They divide my garments among them and cast lots for my clothing''). The passion narrative is told from the start as an exposition of this psalm: suffering as prayer, lament as faith.}

\frage{24. Isaiah 53: Which features of the ``suffering servant of God'' (Isa 53:5.12) recur in our passage?}
\antwort{Isa 53:5 (``crushed for our transgressions'') -- the vicarious suffering indicated by the cry and the death. Isa 53:12 (``he was numbered with the transgressors'') -- literally fulfilled in vv.~27f, two robbers on the right and on the left. Isa 53:7 (silent like a lamb) -- the silent Jesus from v.~20 to v.~37, who speaks only one word.}

\frage{25. Mk 8:31; 9:31; 10:33f: Jesus announces his suffering three times -- which of it is literally fulfilled in vv.~20--27?}
\antwort{The announcements contain the sequence: delivered -- condemned -- handed over to the Gentiles -- mocked -- flogged -- killed. Exactly this sequence runs in 15:1--27 in the same wording (``deliver'', ``condemn'', ``mock'', ``flog'', ``crucify''). Mark shows: the cross is no tragic accident but happens according to Jesus' own understanding of his way -- announced three times, fulfilled ``according to the Scriptures''.}

\frage{26. Mk 10:45: ``\ldots\ to give his life as a ransom for many'' -- how does this word sound in the light of vv.~37.39?}
\antwort{In v.~37 Jesus gives his life: ``Jesus cried out with a loud voice, and gave up the spirit.'' The frame reads: giving of life (v.~37) -- confession (v.~39). What the giving of life effects, the first fruit shows: the enemy of all people comes to faith. The word of Mk 10:45 is narratively redeemed at the cross -- as a word of interpretation, not a formula of calculation.}

\frage{27. Heb 10:19f: What does the torn veil (v.~38) mean in the light of this New Testament interpretation?}
\antwort{Heb 10:19f: ``We have confidence to enter the holy places by the blood of Jesus, by the new and living way that he opened for us through the veil.'' There the veil is an image of Jesus' flesh/death: his death opens the way to God. The tearing ``from top to bottom'' (v.~38) says the same thing narratively: the access stands open -- opened from God, accessible to everyone, no longer through the cult but through Christ.}

\frage{28. Amos 8:9 / Ex 10:21f: Which ancient motif lies behind the darkness (v.~33)?}
\antwort{Amos 8:9: ``On that day I will make the sun go down at noon and darken the earth in broad daylight'' -- an announcement of judgement for the day of the LORD. Ex 10:21f: the darkness as the ninth plague over Egypt -- a sign of judgement. Both together: the darkness from noon to 3 p.m. signals that God's day is breaking in here -- judgement and decision of salvation fall in the same three hours.}

\section*{Closing}

\frage{29. The centurion recognises Jesus in the one who is dying. Where do we encounter people in ``crucified'' situations today?}
\antwort{Open question -- possible impulses for the group: people in illness, loneliness, guilt, despair; the dying whom nobody visits any more; people at the margins. The centurion's gaze can become a measure: \emph{looking} instead of looking away, \emph{recognising} instead of judging. This leads directly into the prayer-request part (7~p.m., 2 groups).}

\vspace{1em}
{\footnotesize Leader's guide · Home group, 8 October 2026 · Bible text: World English Bible (public domain); German edition of the handout: Luther Bible 2017, © Deutsche Bibelgesellschaft.}
\end{document}
